\documentclass[11pt]{article}

\usepackage[a4paper,margin=25mm]{geometry}
\usepackage{amsmath}
\usepackage{amssymb}
\usepackage{graphicx}
\usepackage{hyperref}
\usepackage{siunitx}
\usepackage{xurl}

\title{Rigid Anisotropic Gaussian Model for the 12 CAIN Test Particles}
\author{OPALX Beam--Beam Sandbox}
\date{\today}

\newcommand{\vect}[1]{\boldsymbol{#1}}
\newcommand{\E}{\vect{E}}
\newcommand{\B}{\vect{B}}
\newcommand{\Pmom}{\vect{P}}
\newcommand{\betav}{\boldsymbol{\beta}}

\begin{document}
\maketitle

\section{Purpose}

This note documents the manufactured tracking model implemented in the
\texttt{track12particles.py} sandbox script.  The goal is to compare the
uploaded CAIN test-particle trajectories against a reproducible continuum model
for the beam--beam force.  The model is deliberately simple: the source bunch
is a rigid anisotropic Gaussian moving at constant velocity, with an optional
finite-macroparticle centroid fluctuation.  It is not a PIC model and does not
include bunch disruption, higher-order macro-particle noise, or
self-consistent source evolution.

The CAIN reference file contains twelve right-going test particles:
six electron--positron pairs injected at
\[
  ct/\sigma_z = -1.5,\,-1.0,\,-0.5,\,0,\,0.5,\,1.0,
  \qquad \sigma_z = \SI{0.6}{mm}.
\]
The plotted comparison uses \(s\) in millimetres and \(x\) in micrometres.
Spin columns are parsed but are not part of the present comparison.

\section{Source Selection Convention}

The script supports three source selections:
\begin{itemize}
  \item \texttt{oncoming}: only the electron bunch moving opposite to the
        right-going test particles is included;
  \item \texttt{copropagating}: only the electron bunch moving in the same
        direction as the test particles is included;
  \item \texttt{both}: two equal electron bunches colliding head-on are
        included.
\end{itemize}
The current default for the 12-particle comparison is \texttt{oncoming}.  This
does not mean that both bunches collide head-on.  The head-on two-bunch model is
\texttt{both}.  The \texttt{oncoming} choice is useful because the dominant
force on a right-going ultra-relativistic witness comes from the opposite-going
source; for a co-propagating ultra-relativistic source, the electric and
magnetic parts nearly cancel in the Lorentz force.

\section{Rigid Anisotropic Gaussian Source}

The source is specified by lab-frame beam sizes
\[
  \sigma_x = \sigma_y = \SI{1.944325075701}{\micro m},
  \qquad
  \sigma_z = \SI{0.6}{mm},
\]
and by a source kinetic energy \(K_s=\SI{245}{MeV}\).  The source speed is
\[
  \gamma_s = 1 + \frac{K_s}{m_e c^2},
  \qquad
  \beta_s = \sqrt{1-\gamma_s^{-2}}.
\]
For a source moving along \(\pm z\), the bunch centre is
\[
  \vect r_c(t) = \vect r_c(t_0) + \betav_s c (t-t_0),
  \qquad
  \betav_s = (0,0,\pm\beta_s).
\]
The transverse sizes are unchanged by the boost.  The rest-frame longitudinal
size is dilated:
\[
  \sigma'_x=\sigma_x,\qquad
  \sigma'_y=\sigma_y,\qquad
  \sigma'_z=\gamma_s\sigma_z.
\]

In the source rest frame the charge density is
\[
  \rho'(\vect r')
  =
  \frac{Q}{(2\pi)^{3/2}\sigma'_x\sigma'_y\sigma'_z}
  \exp\left[
    -\frac{1}{2}\left(
      \frac{x'^2}{\sigma_x'^2}
      +\frac{y'^2}{\sigma_y'^2}
      +\frac{z'^2}{\sigma_z'^2}
    \right)
  \right],
\]
with \(Q=-N_e e\), \(N_e=1.25\times 10^{10}\).

\section{Rest-Frame Field Integral}

At a lab event \((t,\vect r)\), the displacement from the moving source centre
is split into transverse and longitudinal components.  For the present
straight-line geometry this is simply
\[
  x'=x-x_c(t),\qquad
  y'=y-y_c(t),\qquad
  z'=\gamma_s\,[z-z_c(t)].
\]
The rest-frame electric field of the triaxial Gaussian is evaluated with the
one-dimensional ellipsoidal integral
\[
  E_i'(\vect r')
  =
  \frac{Q\,x_i'}{4\pi\epsilon_0\sqrt{2\pi}}
  \int_0^\infty
  \frac{
    \exp\left[
      -\frac{1}{2}\sum_j
      \frac{x_j'^2}{\sigma_j'^2+u}
    \right]
  }{
    (\sigma_i'^2+u)
    \sqrt{(\sigma_x'^2+u)(\sigma_y'^2+u)(\sigma_z'^2+u)}
  }\,du .
\]
The script evaluates this integral by trapezoidal quadrature in
\(\log u\).  This is important for the present aspect ratio:
\(\sigma_z'/\sigma_x'\) is large, so a linear quadrature can miss the transverse
variance scale and underestimate the field by orders of magnitude.

\section{Boost to Lab Frame}

The rest-frame magnetic field is zero.  For a boost along \(z\), the lab fields
are
\[
  E_x = \gamma_s E_x',\qquad
  E_y = \gamma_s E_y',\qquad
  E_z = E_z',
\]
and
\[
  \B = \frac{1}{c^2}\,\vect v_s \times \E .
\]
The magnetic field is retained in the particle pusher.  In the present
comparison typical peak fields are in the \(\mathrm{GV/m}\) range and
\(\B\) is of order \(10\)--\(20\,\mathrm{T}\) near overlap.

\section{Witness Push}

The witness momentum is represented in OPALX-normalized units
\[
  \Pmom = \gamma_w \betav_w.
\]
At each substep the script applies the same Boris kick convention used in the
existing manufactured witness prototype:
\begin{align}
  \Pmom^- &=
    \Pmom^n
    + \frac{1}{2}\Delta t\,
      \frac{q c}{m_e c^2}\E^n, \\
  \vect t &=
    \frac{1}{2}\Delta t\,
    \frac{q c^2}{\gamma^- m_e c^2}\B^n,
  \qquad
  \vect s =
    \frac{2\vect t}{1+\vect t\cdot\vect t}, \\
  \Pmom' &= \Pmom^- + \Pmom^- \times \vect t, \\
  \Pmom^+ &= \Pmom^- + \Pmom' \times \vect s, \\
  \Pmom^{n+1} &=
    \Pmom^+
    + \frac{1}{2}\Delta t\,
      \frac{q c}{m_e c^2}\E^n .
\end{align}
The position is advanced with a trapezoidal drift using the velocities before
and after the kick:
\[
  \vect r^{n+1}
  =
  \vect r^n
  +
  \frac{1}{2}
  \left[
    \vect v(\Pmom^n) + \vect v(\Pmom^{n+1})
  \right]\Delta t .
\]

\section{Timestep Handling}

The CAIN file provides output samples on a grid in \(ct\).  The manufactured
tracker uses the same output grid for comparison, but internally substeps each
output interval.  The default maximum substep is
\[
  \Delta t_{\max} = \SI{1}{fs}.
\]
This is intentionally small compared with the overlap transient, and it avoids
confusing output cadence with integration cadence.  The maximum substep can be
changed with the \texttt{max-substep-s} option.

\section{Finite-Macroparticle Centroid Model}

CAIN represents each bunch by a finite random macroparticle sample.  The
manufactured model includes a simple leading-order surrogate for this effect:
the rigid Gaussian source centroid is displaced by a reproducible random offset
with rms
\[
  \sigma_{\bar x_i} = \frac{\sigma_i}{\sqrt{N_{\mathrm{MC}}}},
\]
where \(N_{\mathrm{MC}}\) is the number of source macroparticles.  The default
used here is
\[
  N_{\mathrm{MC}} = 400000.
\]
For the transverse beam size this gives an rms centroid fluctuation of about
\[
  \frac{\SI{1.944}{\micro m}}{\sqrt{400000}}
  \simeq \SI{3.1}{nm}.
\]
The random offset is generated from a fixed seed, so the manufactured trajectory
is reproducible.  This breaks the exact \(y=0\) symmetry and produces a
nonzero vertical field, but it is only a centroid-noise model.  It does not
represent full PIC shot noise, local density fluctuations, or self-consistent
beam disruption.

\section{Current Comparison and Interpretation}

The default command is
\begin{verbatim}
python sandbox/track12particles/track12particles.py \
  --input /Users/adelmann/Desktop/TestParticleOrbit.dat
\end{verbatim}
and writes the CAIN/reference plot, the manufactured-only plot, and the
side-by-side comparison under \texttt{sandbox/track12particles/outputs}.

\begin{figure}[htbp]
  \centering
  \includegraphics[width=\textwidth]{../outputs/track12_cain_vs_manufactured.png}
  \caption{CAIN reference trajectories compared with the rigid anisotropic
  Gaussian manufactured model for \(x(s)\), using the default
  \texttt{oncoming} source selection.}
\end{figure}

\begin{figure}[htbp]
  \centering
  \includegraphics[width=\textwidth]{../outputs/track12_cain_vs_manufactured_y.png}
  \caption{CAIN reference trajectories compared with the rigid anisotropic
  Gaussian manufactured model for \(y(s)\), using the default
  \texttt{oncoming} source selection and the \(1/\sqrt{N_{\mathrm{MC}}}\)
  centroid fluctuation.  The induced manufactured \(y\) motion is much smaller
  than the largest CAIN out-of-plane excursions for the seed used here.}
\end{figure}

\begin{figure}[htbp]
  \centering
  \includegraphics[width=\textwidth]{../outputs/track12_cain_vs_manufactured_z.png}
  \caption{CAIN reference trajectories compared with the rigid anisotropic
  Gaussian manufactured model for longitudinal motion, plotted as \(z=S\)
  versus \(ct\), using the default \texttt{oncoming} source selection.}
\end{figure}

\begin{figure}[htbp]
  \centering
  \includegraphics[width=\textwidth]{../outputs/head_on/track12_cain_vs_manufactured_head_on.png}
  \caption{CAIN reference trajectories compared with the same rigid
  anisotropic Gaussian model for \(x(s)\), using
  \texttt{--source-selection both}, i.e. two equal head-on electron bunches.}
\end{figure}

\begin{figure}[htbp]
  \centering
  \includegraphics[width=\textwidth]{../outputs/head_on/track12_cain_vs_manufactured_y.png}
  \caption{CAIN reference trajectories compared with the rigid anisotropic
  Gaussian manufactured model for \(y(s)\), using
  \texttt{--source-selection both} and the same finite-macroparticle centroid
  fluctuation model.}
\end{figure}

\begin{figure}[htbp]
  \centering
  \includegraphics[width=\textwidth]{../outputs/head_on/track12_cain_vs_manufactured_z.png}
  \caption{CAIN reference trajectories compared with the rigid anisotropic
  Gaussian manufactured model for longitudinal motion, plotted as \(z=S\)
  versus \(ct\), using \texttt{--source-selection both}.}
\end{figure}

\section{OPALX--Manufactured Initial Conditions}

\subsection{Timing Convention from the Source Slides}

The PowerPoint file \texttt{TestParticleOrbitSimulation.pptx} defines the test
particles as artificial insertions during the simulation, not as Breit--Wheeler
outputs.  The relevant timing statement is:
\[
  ct/\sigma_z = -1.5,\,-1.0,\,-0.5,\,0,\,0.5,\,1.0,
\]
with the explicit interpretation that \(ct/\sigma_z=-1.5\) is the moment when
the test particle touches the oncoming bunch edge.  The same slide states that
the test particles are right-going.  Together with the source-slide statement
that the bunch longitudinal Gaussian is cut at \(3\sigma_z\), this fixes the
natural source timing convention.

If the oncoming primary electron bunch centre is at the IP at \(ct=0\), then
for a rigid oncoming source moving toward decreasing \(s\),
\[
  s_c(ct) = -ct .
\]
The initial test-particle coordinates satisfy \(s_i=ct_i\).  Therefore, at the
first timing point \(ct_i=-1.5\sigma_z\), the source centre is at
\(+1.5\sigma_z\), and the test particle is separated from it by
\(-3\sigma_z\), exactly the stated bunch edge.  Thus the primary bunch should
be sampled or evaluated at the test-particle insertion time \(t_i=ct_i/c\), not
at one common \(t=0\) event for all twelve witnesses.

\begin{table}[htbp]
  \centering
  \begin{tabular}{rrrrr}
    \hline
    Pair & \(ct_i/\sigma_z\) & \(ct_i\) [\si{mm}]
    & oncoming \(s_c(ct_i)\) [\si{mm}]
    & \((s_i-s_c)/\sigma_z\) \\
    \hline
    1 & -1.500 & -0.9000 &  0.9000 & -3.000 \\
    2 & -0.999 & -0.5994 &  0.5994 & -1.998 \\
    3 & -0.501 & -0.3006 &  0.3006 & -1.002 \\
    4 &  0.000 &  0.0000 &  0.0000 &  0.000 \\
    5 &  0.501 &  0.3006 & -0.3006 &  1.002 \\
    6 &  0.999 &  0.5994 & -0.5994 &  1.998 \\
    \hline
  \end{tabular}
  \caption{Timing implied by the PowerPoint convention for the six
  right-going test-particle insertion events, using
  \(\sigma_z=\SI{0.6}{mm}\).  The source centre column assumes an oncoming
  primary electron bunch centred at the IP at \(ct=0\).}
  \label{tab:pptx-timing-convention}
\end{table}

This is an important distinction for OPALX comparisons.  A deck that places all
twelve witnesses at their different longitudinal positions at a single OPALX
time \(t=0\), with the primary bunch also centred at the IP at \(t=0\), is a
simultaneous spatial snapshot.  It does not reproduce the PowerPoint insertion
timing for all witnesses except the \(ct=0\) pair.

\subsection{Slide-Timed OPALX Pair Decks}

OPALX has one simulation clock per run.  To reproduce the slide insertion
timing without changing the numerical witness coordinates, the comparison now
uses one electron--positron witness pair per OPALX input deck.  For pair \(i\),
the witness \texttt{FROMFILE} rows retain the TestParticleOrbit values
\[
  (x,y,z,p_x,p_y,p_z)_{\mathrm{OPALX}}
  =
  (x_i,y_i,s_i,P_{x,i}/m_ec,P_{y,i}/m_ec,P_{s,i}/m_ec),
\]
while only the primary source offset is changed:
\[
  R_{0z,\mathrm{source}} = s_{\mathrm{IP}} - ct_i,
  \qquad
  R_{0z,\mathrm{witness}} = s_{\mathrm{IP}}.
\]
Thus at the first OPALX field evaluation,
\[
  s_i - s_c = 2ct_i,
\]
which is the timing convention in Table~\ref{tab:pptx-timing-convention}.
The generated decks are under the OPALX sandbox
\texttt{timing\_pairs/} directory; the generator is
\texttt{generate\_timing\_pair\_inputs.py}.  These decks deliberately fix only
the primary-bunch sampling time.  They do not tune the field normalization,
species mapping, source species, source direction, or pusher details.

\begin{table}[htbp]
  \centering
  \begin{tabular}{rrrrrr}
    \hline
    Pair & \(ct_i/\sigma_z\)
    & \(R_{0z,\mathrm{source}}\) [\si{mm}]
    & \(R_{0z,\mathrm{witness}}\) [\si{mm}]
    & \(s_i-s_c\) [\si{mm}]
    & \((s_i-s_c)/\sigma_z\) \\
    \hline
    1 & -1.500 & 3.9000 & 3.0000 & -1.8000 & -3.000 \\
    2 & -0.999 & 3.5994 & 3.0000 & -1.1988 & -1.998 \\
    3 & -0.501 & 3.3006 & 3.0000 & -0.6012 & -1.002 \\
    4 &  0.000 & 3.0000 & 3.0000 &  0.0000 &  0.000 \\
    5 &  0.501 & 2.6994 & 3.0000 &  0.6012 &  1.002 \\
    6 &  0.999 & 2.4006 & 3.0000 &  1.1988 &  1.998 \\
    \hline
  \end{tabular}
  \caption{Generated OPALX source and witness offsets for the slide-timed
  pair decks.  The metadata check gives
  \(\max |(s_i-s_c)/\sigma_z - 2ct_i/\sigma_z|=0\) to the printed precision.}
  \label{tab:opalx-slide-timed-generated-offsets}
\end{table}

All six slide-timed decks were run with
\(\Delta t=\SI{1}{fs}\), \(N_{\mathrm{source}}=400000\), a
\(32\times 32\times 64\) field mesh, and ten OPALX steps.  The runs used
\texttt{COPY\_TIME=0}, so the witnesses see only the tracked primary source
container.  Every run reached the active BeamBeam state with three active
containers, one source container and one witness in each witness container.
Each run wrote one final H5 group at
\[
  t = \SI{1.0e-14}{s}.
\]
The first-kick CSV diagnostic was enabled with
\texttt{OPALX\_BB\_WITNESS\_KICK\_STEPS=1}.  Each pair CSV contains four data
rows: steps 0 and 1 for the electron and positron witness containers.
The exact OPALX run command used for the six-pair serial validation was:
\begin{footnotesize}
\begin{verbatim}
for pair in 1 2 3 4 5 6; do
  dir="/Users/adelmann/git/opalx-beambeam/sandbox/track12particles/opalx/timing_pairs/pair${pair}"
  stem="track12_pair${pair}_slide_timed_one_source_1fs_400k_10steps"
  echo "running pair ${pair}"
  (cd "$dir" && env OMP_NUM_THREADS=8 \
    OPALX_BB_WITNESS_KICK_CSV="${stem}_kicks.csv" \
    OPALX_BB_WITNESS_KICK_STEPS=1 \
    /Users/adelmann/git/opalx-beambeam/build_openmp/src/opalx "${stem}.in")
done
\end{verbatim}
\end{footnotesize}

\subsection{First-Kick Field Comparison}

Table~\ref{tab:slide-timed-first-kick-fields-oncoming} compares the OPALX
first-kick field magnitudes with the manufactured anisotropic Gaussian at the
same witness positions and at the same slide-timed source offsets.  The table
uses the oncoming manufactured source convention implied by the PowerPoint
timing sketch.

\begin{table}[htbp]
  \centering
  \sisetup{scientific-notation=true, round-mode=figures, round-precision=4}
  \resizebox{\textwidth}{!}{%
  \begin{tabular}{rrrrrr}
    \hline
    Pair & \(s-s_{\mathrm{IP}}\) [\si{mm}]
    & \(|E|_{\mathrm{OPALX}}\) [\si{V/m}]
    & \(|E|_{\mathrm{manuf.}}\) [\si{V/m}]
    & \(|E|_{\mathrm{OPALX}}/|E|_{\mathrm{manuf.}}\)
    & \(|B|_{\mathrm{OPALX}}/|B|_{\mathrm{manuf.}}\) \\
    \hline
    1 & -0.899887 & 3.8787e7 & 5.3768e7 & 0.721 & 0.721 \\
    2 & -0.599287 & 8.4742e8 & 6.5785e8 & 1.288 & 1.288 \\
    3 & -0.300487 & 3.8295e9 & 2.9315e9 & 1.306 & 1.306 \\
    4 &  0.000113 & 6.2682e9 & 4.8444e9 & 1.294 & 1.294 \\
    5 &  0.300713 & 3.8666e9 & 2.9333e9 & 1.318 & 1.318 \\
    6 &  0.599513 & 8.3881e8 & 6.5864e8 & 1.274 & 1.274 \\
    \hline
  \end{tabular}
  }
  \caption{First-kick field magnitudes for the slide-timed pair decks,
  compared with the oncoming manufactured source.  Electron and positron rows
  have the same field magnitudes for each pair, so the table lists the pair
  mean.  \underline{Key result:} the field magnitude disagreement is now at
  the \(0.72\)--\(1.32\) level, not orders of magnitude.}
  \label{tab:slide-timed-first-kick-fields-oncoming}
\end{table}

Repeating the same first-kick field evaluation with a copropagating
manufactured source gives nearly the same field-magnitude ratios:
\[
  0.724 \le
  |E|_{\mathrm{OPALX}}/|E|_{\mathrm{manuf.}}
  \le 1.316,
  \qquad
  0.724 \le
  |B|_{\mathrm{OPALX}}/|B|_{\mathrm{manuf.}}
  \le 1.316.
\]
The equality of the electric and magnetic ratios reflects that both models have
\(|B|\simeq |E|/c\) for this boosted-source geometry.

\subsection{Ten-Step Final-State Comparison}

The final ten-step phase-space comparison is summarized in
Tables~\ref{tab:slide-timed-final-oncoming} and
\ref{tab:slide-timed-final-copropagating}.  Both tables use the H5 charge
assignment for the OPALX witness containers: container 1 is the electron
witness and container 2 is the positron witness.  The reported values are the
maximum absolute OPALX-minus-manufactured differences over the two witness
species for each pair.

\begin{table}[htbp]
  \centering
  \begin{tabular}{rrrr}
    \hline
    Pair & \(\max|\Delta x|\) [\si{\micro m}]
    & \(\max|\Delta p_x|\)
    & \(\max|\Delta s|\) [\si{\micro m}] \\
    \hline
    1 & 0.000467 (0,0\%) & 0.000626 (105,1\%) & 0.000000051 (0,0\%) \\
    2 & 0.005925 (0,3\%) & 0.007929 (109,2\%) & 0.000006665 (0,0\%) \\
    3 & 0.026389 (1,3\%) & 0.035290 (109,3\%) & 0.000131402 (0,0\%) \\
    4 & 0.043494 (2,2\%) & 0.058121 (109,3\%) & 0.000357354 (0,0\%) \\
    5 & 0.026306 (1,3\%) & 0.035082 (109,5\%) & 0.000130204 (0,0\%) \\
    6 & 0.005873 (0,3\%) & 0.007815 (109,2\%) & 0.000006535 (0,0\%) \\
    \hline
  \end{tabular}
  \caption{Ten-step final-state differences for the slide-timed OPALX pair
  decks compared with an oncoming manufactured source.  Percentages in
  parentheses are relative to the corresponding manufactured value and are
  rounded to one digit after the decimal comma.  The largest differences
  are \(\max|\Delta x|=\SI{0.04349}{\micro m}\),
  \(\max|\Delta p_x|=0.05812\), and
  \(\max|\Delta s|=\SI{3.57e-4}{\micro m}\).}
  \label{tab:slide-timed-final-oncoming}
\end{table}

\begin{table}[htbp]
  \centering
  \begin{tabular}{rrrr}
    \hline
    Pair & \(\max|\Delta x|\) [\si{\micro m}]
    & \(\max|\Delta p_x|\)
    & \(\max|\Delta s|\) [\si{\micro m}] \\
    \hline
    1 & 0.000054 & 0.000072 & 0.000000006 \\
    2 & 0.000885 & 0.001181 & 0.000000049 \\
    3 & 0.003981 & 0.005311 & 0.000000528 \\
    4 & 0.006549 & 0.008740 & 0.000001231 \\
    5 & 0.004011 & 0.005352 & 0.000000544 \\
    6 & 0.000884 & 0.001179 & 0.000000045 \\
    \hline
  \end{tabular}
  \caption{Ten-step final-state differences for the same slide-timed OPALX
  pair decks compared with a copropagating manufactured source.  The largest
  differences are \(\max|\Delta x|=\SI{0.00655}{\micro m}\),
  \(\max|\Delta p_x|=0.00874\), and
  \(\max|\Delta s|=\SI{1.23e-6}{\micro m}\).}
  \label{tab:slide-timed-final-copropagating}
\end{table}

\noindent\underline{Numerical conclusion:} after fixing the slide timing, the
current OPALX one-source decks agree much more closely with a copropagating
manufactured source than with an oncoming manufactured source.  This is a
convention result, not a field or pusher tuning result.  The
slide-timed OPALX decks still use the existing tracked primary source momentum
with positive \(p_z\).  Therefore the next physics decision is to explicitly
choose the OPALX primary-source direction convention for the intended oncoming
comparison before changing field normalization, species mapping, or Boris
pusher details.

For the early one-source diagnostic, CAIN was not used as an active reference
model.  The same twelve numerical witness-particle initial conditions were
read into OPALX through two \texttt{FROMFILE} distributions and into the
manufactured Python model.  Table~\ref{tab:opalx-manufactured-initial-conditions}
compares the OPALX input files against the manufactured-script initial state.
The reported deltas are OPALX minus manufactured.  The agreement is exact for
all listed coordinates and momenta.  The exact-timed comparison described
below restores the full CAIN trajectories as the pointwise reference.

\begin{table}[htbp]
  \centering
  \sisetup{
    round-mode=places,
    round-precision=4,
    table-number-alignment=center
  }
  \resizebox{\textwidth}{!}{%
  \begin{tabular}{llrrrrrrrrrrrr}
    \hline
    Species & Pair
    & \(x\) [\si{\micro m}]
    & \(y\) [\si{\micro m}]
    & \(s\) [\si{mm}]
    & \(p_x\)
    & \(p_y\)
    & \(p_z\)
    & \(\Delta x\) [\si{\micro m}]
    & \(\Delta y\) [\si{\micro m}]
    & \(\Delta s\) [\si{\micro m}]
    & \(\Delta p_x\)
    & \(\Delta p_y\)
    & \(\Delta p_z\) \\
    \hline
    electron & 1 & 1.9443250757 & 0 & -0.9000 & 0 & 0 & 1.7320511804 & 0 & 0 & 0 & 0 & 0 & 0 \\
    electron & 2 & 1.9443250757 & 0 & -0.5994 & 0 & 0 & 1.7320511804 & 0 & 0 & 0 & 0 & 0 & 0 \\
    electron & 3 & 1.9443250757 & 0 & -0.3006 & 0 & 0 & 1.7320511804 & 0 & 0 & 0 & 0 & 0 & 0 \\
    electron & 4 & 1.9443250757 & 0 &  0.0000 & 0 & 0 & 1.7320511804 & 0 & 0 & 0 & 0 & 0 & 0 \\
    electron & 5 & 1.9443250757 & 0 &  0.3006 & 0 & 0 & 1.7320511804 & 0 & 0 & 0 & 0 & 0 & 0 \\
    electron & 6 & 1.9443250757 & 0 &  0.5994 & 0 & 0 & 1.7320511804 & 0 & 0 & 0 & 0 & 0 & 0 \\
    positron & 1 & 1.9443250757 & 0 & -0.9000 & 0 & 0 & 1.7320511804 & 0 & 0 & 0 & 0 & 0 & 0 \\
    positron & 2 & 1.9443250757 & 0 & -0.5994 & 0 & 0 & 1.7320511804 & 0 & 0 & 0 & 0 & 0 & 0 \\
    positron & 3 & 1.9443250757 & 0 & -0.3006 & 0 & 0 & 1.7320511804 & 0 & 0 & 0 & 0 & 0 & 0 \\
    positron & 4 & 1.9443250757 & 0 &  0.0000 & 0 & 0 & 1.7320511804 & 0 & 0 & 0 & 0 & 0 & 0 \\
    positron & 5 & 1.9443250757 & 0 &  0.3006 & 0 & 0 & 1.7320511804 & 0 & 0 & 0 & 0 & 0 & 0 \\
    positron & 6 & 1.9443250757 & 0 &  0.5994 & 0 & 0 & 1.7320511804 & 0 & 0 & 0 & 0 & 0 & 0 \\
    \hline
  \end{tabular}%
  }
  \caption{Initial witness-particle phase space used by both OPALX and the
  manufactured Python model.  The momenta are normalized OPALX momenta
  \(p_i=\beta_i\gamma\).  All deltas are OPALX minus manufactured.}
  \label{tab:opalx-manufactured-initial-conditions}
\end{table}

\section{Exact-Timed OPALX--CAIN Comparison}

The current workflow is implemented in
\texttt{sandbox/track12particles/opalx/timed}.  It replaces the six independent
pair decks by one physical clock, one BeamBeam element, and two emitted witness
containers.  Each emitted-file record uses the named format
\begin{verbatim}
x y z px py pz birth_time
\end{verbatim}
where position is relative to the emission-source reference position,
momentum is normalized by \(m_e c\), and \texttt{birth\_time} is an exact time
offset from the emission-source \texttt{T0}.  Electron and positron files have
identical numeric rows; the containing OPALX beam supplies the species and
charge sign.

The CAIN output interval is
\[
  \Delta(ct)=\SI{1.8}{\micro m},\qquad
  \Delta t=\frac{\Delta(ct)}{c}
           =\SI{6.004153713566736}{fs}.
\]
The tracker begins one interval before the first birth so that the mirrored
primary source exists before a witness is emitted, and the final global step
is 1500.  Hence every OPALX sample lies on a CAIN sample time and no temporal
interpolation is performed.  Floating-point comparison at an exact step
boundary can make the first H5 record be reference step zero or one.  The exact
input row is used as sample zero when H5 begins at step one.

The physical primary and its reflection about the midpoint of the BeamBeam
element supply the co-propagating and oncoming fields.  The witness containers
are excluded from deposition because twelve physical test particles are far
too few to represent a charge density, but they gather both source fields.
The source sample is held rigid with \texttt{BBRIGID=TRUE}.  The fixed field
domain is a \(\SI{2.4}{mm}\times\SI{0.24}{mm}\) transverse rectangle.
\texttt{DELETEONTRANSVERSEEXIT=FALSE} prevents deletion, but the particle
layout remains periodic; any boundary crossing must therefore be diagnosed and
cannot be interpreted as an open-domain orbit.

The production discretization uses 400000 deterministic primary
macroparticles and a \(1024\times128\times128\) mesh.  The nominal transverse
cell sizes are approximately \(\SI{2.34}{\micro m}\) and
\(\SI{1.88}{\micro m}\).  Coarser local cases validate parsing, emission,
container assignment, and output matching only; they do not resolve the
\(\sigma_x=\sigma_y=\SI{1.9443}{\micro m}\) primary density and must not be
interpreted as a physics comparison.

The A100 executable is built for compute capability 8.0 with Kokkos 5.2.0.
The production run uses four MPI ranks and four A100 GPUs, with one device per
rank selected by \texttt{--kokkos-map-device-id-by=mpi\_rank}.  The two witness
H5 files are copied back after completion and analyzed locally.  Stored H5
particle coordinates are offsets from the three-component \texttt{RefPartR}
attribute; all three reference components are added when reconstructing
absolute trajectories.

The four-rank, four-A100 production run completed as Merlin job 353971 in
1:43:15.  All twelve witnesses were emitted in the intended order, no witness
charge entered the field solve, and all 13012 CAIN rows were matched by
species, pair, and exact grid time.  Raw H5 particle IDs are not persistent
under MPI redistribution: the electron file contains 10 distinct raw IDs and
the positron file 16 for six physical trajectories per species.  The analysis
therefore assigns pair identity by birth order and one-step phase-space
continuity.

Table~\ref{tab:track12-first-kicks} records the first nonzero transverse
momentum sample from the original, pre-reference-offset-fix run.  The positron
values are equal and opposite to the electron values to the printed precision,
confirming that the two witness containers have the intended charge signs.
Only pair 1 has a comparable kick; pairs 3 and 6 have the wrong sign, and the
central pair receives only 1.20\% of the CAIN kick.

\begin{table}[htbp]
  \centering
  \begin{tabular}{rrrr}
    \hline
    Pair & $p_{x,\mathrm{CAIN}}$ & $p_{x,\mathrm{OPALX}}$
         & OPALX/CAIN \\
    \hline
    1 & 2.2302e-5 & 3.0190e-5  &  1.3537 \\
    2 & 3.8885e-4 & 1.6348e-5  &  0.0420 \\
    3 & 6.6503e-3 & -5.8373e-5 & -0.0088 \\
    4 & 2.2261e-2 & 2.6816e-4  &  0.0120 \\
    5 & 1.2516e-2 & 4.8192e-5  &  0.0039 \\
    6 & 1.7218e-3 & -6.2312e-6 & -0.0036 \\
    \hline
  \end{tabular}
  \caption{Pre-fix first electron transverse-momentum sample on the exact CAIN
  grid.  The positron samples have the opposite sign.}
  \label{tab:track12-first-kicks}
\end{table}

The full pointwise error is summarized in
Table~\ref{tab:track12-production-errors}.  The small absolute positron error
must not be read as agreement: its relative $L_2$ error is 0.943 because the
CAIN positron excursions themselves are small.

\begin{table}[htbp]
  \centering
  \begin{tabular}{lrrrrr}
    \hline
    Species & Samples & RMSE $x$ [\si{\micro m}] & relative $L_2(x)$
            & RMSE $y$ [\si{\micro m}] & RMSE $s$ [\si{\micro m}] \\
    \hline
    electron & 6506 & 511.68 & 1.12695 & 251.14 & 30.52 \\
    positron & 6506 &   5.04 & 0.94299 &   1.08 &  0.78 \\
    \hline
  \end{tabular}
  \caption{Pre-fix four-A100 OPALX trajectory error relative to all CAIN
  samples.}
  \label{tab:track12-production-errors}
\end{table}

The fixed particle layout is periodic transversely even though the Poisson
boundary condition is open.  Four over-kicked electron trajectories cross the
$y$ boundary, for nine wrap events in total.  The analysis unwraps those H5
coordinates to preserve continuity, but after the first wrap the simulated
particle samples the field on the opposite side of the numerical domain and
is no longer a physical open-boundary trajectory.

\subsection{Witness Reference-Frame Repair and Field Convergence}

The original trajectory failure was subsequently traced to the conversion
between independently referenced particle containers.  Particle position
\texttt{R} is stored relative to its container's reference particle.  For the
track12 witnesses, \texttt{RefPartR.x}$=\sigma_x$ while particle-local
\texttt{R.x} is nearly zero.  The BeamBeam witness gather translated only the
longitudinal path offset, so the pre-fix run sampled the field near the source
axis even though the reconstructed H5 position correctly appeared at
$x=\sigma_x$.

The corrected gather includes the transverse reference-particle difference
\[
  \vect R_{\mathrm{ref},w}-\vect R_{\mathrm{ref},s}
\]
when mapping witness-local positions to the source frame, while retaining the
curvilinear path difference $s_w-s_s$ for the longitudinal component.  The
inverse transform and field rotation are unchanged.  Thus this repair changes
only the point at which a passive witness samples the already-computed field;
it does not alter deposition, charge normalization, the Poisson solve, field
units, or the pusher.

The persistent pair-4 diagnostic compares the field at $x=\sigma_x$ against
the untruncated rigid-Gaussian solution.  Table~\ref{tab:track12-field-convergence}
shows monotonic refinement for both the normal seeded source and a deterministic
equal-probability tensor cubature.  The listed domain is the full transverse
width in each direction.

\begin{table}[htbp]
  \centering
  \begin{tabular}{rrrr}
    \hline
    Mesh & Cell [\si{\micro m}] & Seeded $|E|$ ratio & Tensor $|E|$ ratio \\
    \hline
    $32\times32$   & 1.290 & 0.901839 & 0.893207 \\
    $64\times64$   & 0.635 & 0.958711 & 0.953227 \\
    $128\times128$ & 0.315 & 0.984256 & 0.978678 \\
    \hline
  \end{tabular}
  \caption{OPALX-to-analytic electric-field magnitude ratio on a full
  \SI{40}{\micro m}$\times$\SI{40}{\micro m} domain after the reference-frame
  repair.  The scan case name \texttt{square20} denotes the
  \SI{20}{\micro m} half-width.}
  \label{tab:track12-field-convergence}
\end{table}

At approximately \SI{0.63}{\micro m} cell spacing, doubling the full domain
from \SI{40}{\micro m} to \SI{80}{\micro m} or changing to the 2:1 aspect
ratio changes the tensor result by less than 0.1\%.  Seeded and tensor results
differ by approximately 1\% or less.
The earlier apparent domain/rank pathology was therefore dominated by the
missing transverse reference offset and by treating \texttt{CIRCLE} aperture
arguments as radii even though they are full diameters.

The exact fixed-source production probe used 400000 particles, a
$1024\times128\times128$ mesh, and the
\SI{2.4}{mm}$\times$\SI{0.24}{mm} domain.  Four-rank Merlin job 353974
completed on four A100s in \SI{23.14}{s}.  Its central pair result is
summarized in Table~\ref{tab:track12-corrected-pair4-probe}.

\begin{table}[htbp]
  \centering
  \begin{tabular}{lr}
    \hline
    Quantity & Value \\
    \hline
    OPALX $|E|$ [\si{V/m}] & $6.569905\times10^9$ \\
    OPALX/analytic $|E|$ & 0.678097 \\
    Electric-field direction cosine & 0.999995 \\
    Field-predicted OPALX $\Delta p_x$ & 0.023143 \\
    CAIN first $\Delta p_x$ & 0.022261 \\
    OPALX/CAIN first kick & 1.039622 \\
    \hline
  \end{tabular}
  \caption{Corrected pair-4 field probe on the exact production
  discretization.  The analytic denominator is an untruncated continuum
  Gaussian, while OPALX uses the fixed finite source and production mesh.}
  \label{tab:track12-corrected-pair4-probe}
\end{table}

The corrected production-discretization field therefore predicts the central
CAIN first kick within 3.96\%, instead of the pre-fix trajectory's 1.20\%
residual.  The timed-input, species, and comparison pipeline remains validated,
and a complete corrected 12-particle production rerun is now justified.  That
rerun has not yet been performed; the trajectory tables above remain pre-fix
baselines.

\end{document}
